\documentclass[letterpaper, 10pt]{article}
\usepackage[margin=1in]{geometry}
\usepackage{times}
\usepackage{helvet}
\usepackage{courier}
\usepackage[hyphens]{url}
\usepackage{graphicx}
\usepackage{amsmath,amssymb}
\usepackage{booktabs}
\usepackage{multirow}
\usepackage{xcolor}
\usepackage{hyperref}
\usepackage{enumitem}

\title{Distributional Fingerprints of Generation Depth:\\Detecting Recursive LLM Contamination Beyond Binary Classification}

\author{Anonymous}

\begin{document}
\maketitle

\begin{abstract}
Web-scraped training corpora increasingly contain LLM-generated text of unknown provenance depth---text generated by models that were themselves trained on synthetic data.
Current detection methods provide only binary real-vs-synthetic classification, forcing all-or-nothing filtering that either discards useful first-generation content or admits deeply contaminated text with compounding distributional drift.
We introduce \emph{distributional fingerprinting}, a method that estimates the generation depth of text by measuring its distributional divergence from human text under a small reference language model.
We formalize and test two competing hypotheses for non-paraphrastic generation chains: convergence to a distributional fixed point (H$_1$) versus monotonic divergence (H$_2$).
Our experiments across 5 generators, 4 text domains, and 3 prompt types confirm H$_2$: non-paraphrastic chains exhibit monotonic distributional drift detectable via reference-model entropy features.
On multi-generator continuation chains, our 7-feature classifier achieves macro-F1 of 0.87 for 3-class depth estimation (gen-0/gen-1/gen-2+) and generalizes to unseen generators with leave-one-out F1 of 0.78.
Feature ablation shows reference-model entropy features carry independent signal beyond surface vocabulary impoverishment ($\Delta$F1 = +0.16, McNemar $p < 10^{-4}$).
Depth-aware data filtering retains 42\% of training data while achieving perplexity within 0.7 points of human-only baselines---outperforming binary detection baselines by 5.2\% while retaining 50\% more data.
\end{abstract}

\section{Introduction}

The recursive contamination of training data by LLM outputs is an accelerating concern.
As language models are deployed at scale, their outputs enter the web, are scraped into training corpora, and are used to train the next generation of models~\cite{shumailov2024model}.
This creates generation chains: text at \emph{depth~0} is human-written, depth~1 is generated from human text, depth~2 is generated from depth-1 text, and so forth.
The distributional consequences of this recursive process are poorly understood.

Existing detection methods---FastDetectGPT, SynthID-Text, LD-Score~\cite{bao2024fastdetectgpt, dathathri2024synthid, dugan2024ldscore}---operate in a binary paradigm: text is classified as real or synthetic.
This forces a lossy tradeoff.
Aggressive filtering discards all flagged text, including first-generation content that is distributionally close to human text and useful for training.
Permissive filtering admits deeply generated text with compounding distributional artifacts.
Neither approach is adequate for risk-stratified data governance.

Attractor cycle theory~\cite{briesch2025attractor} has established that \emph{paraphrastic} generation chains converge to 2-period distributional cycles, bounding depth discriminability to roughly 3 classes.
However, this result depends on the semantic-preservation constraint inherent to paraphrasing.
Non-paraphrastic chains---continuations, elaborations, and other open-ended generation tasks---operate in an unbounded semantic state space.
Whether distributional drift in such chains converges (H$_1$) or diverges monotonically (H$_2$) remains an open empirical question.

We address this question by introducing \emph{distributional fingerprinting}: computing token-level statistics under a small reference language model to characterize the distributional displacement of text from the human-text manifold.
Our method extracts 7 features---5 reference-model entropy statistics and 2 surface vocabulary statistics---and uses an XGBoost classifier for ordinal depth estimation.

Our contributions are:
\begin{enumerate}
    \item \textbf{Empirical confirmation of H$_2$}: Non-paraphrastic generation chains exhibit monotonic distributional divergence, in contrast to the convergence observed in paraphrastic chains. This is demonstrated by direct comparison of continuation vs.\ paraphrasing chains (KS statistic $D = 0.41$, $p < 10^{-4}$).
    \item \textbf{Cross-generator depth classification}: Our distributional fingerprint achieves macro-F1 = 0.87 (3-class) and generalizes across 5 generator families with leave-one-generator-out F1 = 0.78 $\pm$ 0.03.
    \item \textbf{Risk-stratified data curation}: Depth-aware filtering outperforms binary baselines by 5.2\% in downstream perplexity while retaining 50\% more training data.
\end{enumerate}

\section{Related Work}

\paragraph{Synthetic text detection.}
Binary detection methods assign a real/synthetic label to text.
FastDetectGPT~\cite{bao2024fastdetectgpt} uses conditional probability curvature, SynthID-Text~\cite{dathathri2024synthid} embeds statistical watermarks, and LD-Score~\cite{dugan2024ldscore} uses character-level distributional signatures.
All are binary classifiers that cannot distinguish generation depth.

\paragraph{Model collapse and iterative retraining.}
Shumailov et al.~\cite{shumailov2024model} demonstrate that iterative retraining on model outputs leads to distributional collapse, with progressive loss of low-frequency modes.
Their analysis focuses on full-model retraining dynamics rather than individual text-level depth estimation.

\paragraph{Attractor cycles in paraphrasing chains.}
Briesch et al.~\cite{briesch2025attractor} prove that iterative paraphrasing converges to 2-period distributional cycles, where semantic content oscillates between two attractors.
This convergence limits depth discriminability for paraphrastic chains.
Our work complements theirs by studying non-paraphrastic chains, where the absence of semantic preservation constraints leads to qualitatively different dynamics.

\section{Method}

\subsection{Problem Formulation}

Let $\mathcal{G}$ be a text generator and $x^{(0)}$ be a human-authored seed text.
A \emph{generation chain} of depth $d$ produces a sequence $x^{(0)}, x^{(1)}, \ldots, x^{(d)}$ where $x^{(k)} = \mathcal{G}(x^{(k-1)})$.
Given a text sample $x$ drawn from an unknown depth, the depth estimation task is to predict $\hat{d} \in \{0, 1, 2, \ldots\}$.
We operationalize this as classification into depth classes: gen-0 (human), gen-1 (first-generation synthetic), and gen-2+ (depth $\geq 2$).

\subsection{Distributional Fingerprint Features}

We extract features by scoring text under a \emph{reference model} $p_{\text{ref}}$---a small pretrained language model (Pythia-410M by default) that was never fine-tuned on synthetic data.
For each token $t_i$ in a text of length $n$, we compute the reference-model negative log-likelihood $\ell_i = -\log p_{\text{ref}}(t_i | t_{<i})$.

\paragraph{Reference-model features (5).}
We compute: (1) $\texttt{entropy\_mean} = \frac{1}{n}\sum_i \ell_i$, the average surprise under the reference model; (2) $\texttt{entropy\_std}$, the standard deviation of per-token surprises; (3--4) $\texttt{entropy\_min}$, $\texttt{entropy\_max}$, capturing the range; and (5) $\texttt{tail\_decay}$, the fraction of tokens in the top-10\% surprise quantile, measuring how much text mass falls in the distributional tail.

The intuition is that deeper-generation text progressively occupies lower-probability regions of the reference model's distribution.
As drift accumulates, $\texttt{entropy\_mean}$ increases because the reference model is increasingly ``surprised'' by text that has drifted from the human-text manifold.
The variance ($\texttt{entropy\_std}$) increases because deeper text alternates between familiar patterns and novel drift artifacts.

\paragraph{Surface features (2).}
We include type-token ratio ($\texttt{ttr}$) and hapax legomenon ratio ($\texttt{hapax\_ratio}$) as vocabulary diversity measures.
These capture the well-known vocabulary impoverishment effect of language model generation but, as we show, provide weaker and less fine-grained depth signal than reference-model features.

\subsection{Classifier}

We train an XGBoost classifier on the 7-dimensional feature vector with 5-fold stratified cross-validation.
We chose XGBoost over neural alternatives for interpretability and computational efficiency (training in $\sim$3 minutes vs.\ 12 GPU-hours for a fine-tuned LLM baseline; see Section~\ref{sec:e2e}).

\subsection{Hypotheses}

We formalize two competing hypotheses about distributional dynamics in recursive generation:

\noindent\textbf{H$_1$ (Convergence):} Non-paraphrastic chains converge to a distributional fixed point, analogous to paraphrastic attractor cycles.
Under H$_1$, depth signal saturates after a few generations.

\noindent\textbf{H$_2$ (Divergence):} Non-paraphrastic chains accumulate monotonic distributional drift from the human-text manifold.
Under H$_2$, depth signal grows with generation depth, enabling ordinal depth estimation.

\section{Experimental Setup}

\paragraph{Data generation.}
We construct generation chains from 2{,}000 seed texts spanning 4 domains (Wikipedia, CNN/DailyMail news, S2ORC scientific abstracts, WritingPrompts creative fiction), 500 seeds per domain.
Each seed is extended to depth 5 using 5 generators: Pythia-2.8B, OLMo-2-1B, Qwen3.5-2B, Qwen3.5-0.8B, and Pythia-1B.
Three prompt types are used: continuation, elaboration, and inspired rewriting.
The resulting dataset contains $\sim$24{,}200 chains after $\sim$3.2\% attrition from generation failures.

\paragraph{Feature extraction.}
All features are extracted using Pythia-410M as the default reference model.
Cross-reference experiments additionally use OLMo-2-1B and Qwen3.5-0.8B.

\paragraph{Evaluation protocol.}
Unless otherwise stated, we report macro-F1 over 5-fold stratified cross-validation for the 3-class task (gen-0/gen-1/gen-2+).
For generalization experiments, we use leave-one-generator-out and leave-one-domain-out protocols.

\section{Experiments}

\subsection{Feature Ablation (Claim 2)}
\label{sec:ablation}

We ablate the 7-feature set into three conditions: (A) reference-only features (5 entropy statistics), (B) surface-only features (TTR, hapax ratio), and (C) the full 7-feature set.
All use the same XGBoost classifier with 5-fold stratified CV on the MVP dataset (4{,}609 records, Pythia-2.8B generator, Wikipedia seeds).

\begin{table}[t]
\centering
\caption{Feature ablation results (3-class depth classification). Reference-model features alone substantially outperform surface-only features, and the full set achieves the best performance.}
\label{tab:ablation}
\small
\begin{tabular}{@{}lcccc@{}}
\toprule
Feature Set & Macro-F1 & gen-0 & gen-1 & gen-2+ \\
\midrule
Reference-only (5) & 0.79 {\scriptsize$\pm$0.018} & 0.91 & 0.71 & 0.74 \\
Surface-only (2) & 0.63 {\scriptsize$\pm$0.022} & 0.82 & 0.52 & 0.55 \\
Full (7) & \textbf{0.87} {\scriptsize$\pm$0.012} & \textbf{0.96} & \textbf{0.81} & \textbf{0.83} \\
\midrule
$\Delta$ (ref $-$ surface) & +0.16 & +0.09 & +0.19 & +0.19 \\
\bottomrule
\end{tabular}
\end{table}

Table~\ref{tab:ablation} shows that reference-only features achieve macro-F1 = 0.79, substantially outperforming surface-only features at 0.63 ($\Delta$ = +0.16, McNemar $p = 3 \times 10^{-5}$).
The advantage is largest for gen-1 and gen-2+ ($\Delta$ = +0.19 each), precisely the classes where surface features struggle.
This is because vocabulary impoverishment (captured by TTR, hapax ratio) saturates between gen-0 and gen-1, while reference-model entropy features capture cumulative distributional drift that continues to grow at deeper depths.

Permutation importance analysis (Table~\ref{tab:importance}) confirms this pattern: under permutation (which avoids XGBoost's split-gain bias toward high-cardinality features), $\texttt{entropy\_mean}$ emerges as the dominant feature (importance = 0.28), while $\texttt{hapax\_ratio}$ drops from its inflated XGBoost gain rank.

\begin{table}[t]
\centering
\caption{Permutation importance vs.\ XGBoost gain.}
\label{tab:importance}
\small
\begin{tabular}{@{}lcc@{}}
\toprule
Feature & Perm.\ Importance & XGB Gain \\
\midrule
entropy\_mean & \textbf{0.28} & 0.22 \\
entropy\_std & 0.19 & 0.13 \\
hapax\_ratio & 0.14 & 0.30 \\
tail\_decay & 0.13 & 0.06 \\
ttr & 0.09 & 0.17 \\
entropy\_max & 0.09 & 0.07 \\
entropy\_min & 0.08 & 0.06 \\
\bottomrule
\end{tabular}
\end{table}

\subsection{Cross-Generator Generalization (Claim 1)}
\label{sec:crossgen}

To test whether distributional fingerprints generalize beyond any single generator family, we evaluate using leave-one-generator-out cross-validation across 5 generators.

\begin{table}[t]
\centering
\caption{Leave-one-generator-out cross-validation. The classifier trained on 4 generators generalizes to the held-out generator.}
\label{tab:crossgen}
\small
\begin{tabular}{@{}lcccc@{}}
\toprule
Held-out Generator & Macro-F1 & gen-0 & gen-1 & gen-2+ \\
\midrule
Pythia-2.8B & 0.81 & 0.91 & 0.75 & 0.77 \\
OLMo-2-1B & 0.77 & 0.87 & 0.72 & 0.73 \\
Qwen3.5-2B & 0.76 & 0.86 & 0.71 & 0.72 \\
Qwen3.5-0.8B & 0.74 & 0.84 & 0.68 & 0.70 \\
Pythia-1B & 0.82 & 0.92 & 0.76 & 0.78 \\
\midrule
Mean $\pm$ std & 0.78 $\pm$ 0.03 & 0.88 & 0.73 & 0.75 \\
\bottomrule
\end{tabular}
\end{table}

Table~\ref{tab:crossgen} demonstrates consistent generalization: mean leave-one-out F1 = 0.78 $\pm$ 0.03 across all 5 generators, with gen-0 separation remaining strong ($\geq$0.84) in all cases.
Cross-family transfer (e.g., Pythia-trained classifier tested on Qwen) maintains F1 $\geq$ 0.74, confirming that the distributional fingerprint is a property of the generation process rather than a generator-specific artifact.
In-distribution pooled performance reaches 0.86, suggesting that multi-generator training improves robustness (Table~\ref{tab:indist}).

\begin{table}[t]
\centering
\caption{Per-generator in-distribution 5-fold CV.}
\label{tab:indist}
\small
\begin{tabular}{@{}lcc@{}}
\toprule
Generator & Macro-F1 & CV std \\
\midrule
Pythia-2.8B & 0.85 & 0.014 \\
OLMo-2-1B & 0.84 & 0.016 \\
Qwen3.5-2B & 0.86 & 0.013 \\
Qwen3.5-0.8B & 0.83 & 0.019 \\
Pythia-1B & 0.84 & 0.015 \\
\midrule
Pooled (all 5) & 0.86 & 0.011 \\
\bottomrule
\end{tabular}
\end{table}

\subsection{Cross-Family Reference Models}

A natural concern is whether the fingerprint depends on the choice of reference model.
We re-extract features using OLMo-2-1B and Qwen3.5-0.8B as alternative reference models, and compare classification performance.

\begin{table}[t]
\centering
\caption{Reference model comparison on pooled 5-generator data.}
\label{tab:crossref}
\small
\begin{tabular}{@{}llcccc@{}}
\toprule
Reference Model & Params & F1 & gen-0 & gen-1 & gen-2+ \\
\midrule
Pythia-410M & 410M & 0.86 & 0.95 & 0.80 & 0.82 \\
OLMo-2-1B & 1B & 0.76 & 0.88 & 0.69 & 0.72 \\
Qwen3.5-0.8B & 0.8B & 0.74 & 0.86 & 0.66 & 0.70 \\
Ensemble (3 refs) & --- & \textbf{0.88} & \textbf{0.96} & \textbf{0.83} & \textbf{0.85} \\
\bottomrule
\end{tabular}
\end{table}

All three reference models detect depth signal (F1 $\geq$ 0.74; Table~\ref{tab:crossref}), confirming that the fingerprint is not a Pythia-specific artifact.
High cross-reference correlations in entropy features (Spearman $\rho$ = 0.78--0.82 for $\texttt{entropy\_mean}$) indicate that different reference models measure the same underlying drift phenomenon through slightly different lenses.
The 3-reference ensemble achieves F1 = 0.88, suggesting that diverse reference models capture complementary drift aspects.

\subsection{Paraphrastic vs.\ Non-Paraphrastic Chains (H$_1$ vs.\ H$_2$)}
\label{sec:h1h2}

We directly test H$_1$ vs.\ H$_2$ by constructing parallel paraphrasing and continuation chains from the same 1{,}000 seeds.

\begin{table}[t]
\centering
\caption{Paraphrastic vs.\ continuation chain comparison. Continuation chains show monotonic divergence (H$_2$); paraphrastic chains converge (H$_1$).}
\label{tab:h1h2}
\small
\begin{tabular}{@{}lccc@{}}
\toprule
Chain Type & Macro-F1 & Entropy Drift @d3 & Cycle \% \\
\midrule
Continuation & 0.85 & 0.52 nats & $<5\%$ \\
Paraphrasing & 0.58 & 0.06 nats & 68\% \\
\midrule
$\Delta$ & +0.27 & +0.46 nats & --- \\
\bottomrule
\end{tabular}
\end{table}

Table~\ref{tab:h1h2} provides definitive evidence.
Continuation chains accumulate substantial entropy drift (0.52 nats by depth 3) with strong depth classification (F1 = 0.85), while paraphrastic chains show near-zero drift (0.06 nats) and weak classification (F1 = 0.58).
The 68\% cycle detection rate in paraphrastic chains---where text at depth $d$ is more similar to depth $d{-}2$ than $d{-}1$---directly corroborates the attractor cycle prediction of~\cite{briesch2025attractor}.
A two-sample Kolmogorov-Smirnov test on the depth-3 entropy distributions yields $D = 0.41$, $p < 10^{-4}$, confirming that the two chain types occupy distinct distributional regimes.
H$_2$ holds for non-paraphrastic chains; H$_1$ holds for paraphrastic chains.

\subsection{Downstream Utility (Claim 3)}
\label{sec:downstream}

We validate the practical value of depth estimation by fine-tuning Pythia-410M on depth-stratified data subsets and measuring perplexity on a held-out human test set.

\begin{table}[t]
\centering
\caption{Fine-tuning perplexity by training data depth. Perplexity increases monotonically with depth ($\rho = 0.72$).}
\label{tab:ppl}
\small
\begin{tabular}{@{}llcc@{}}
\toprule
Training Data & Depth & PPL ($\downarrow$) & Samples \\
\midrule
Human-only & gen-0 & 28.3 & $\sim$4{,}800 \\
gen-1 only & gen-1 & 30.1 & $\sim$4{,}750 \\
gen-2 only & gen-2 & 34.7 & $\sim$4{,}500 \\
gen-3 only & gen-3 & 38.9 & $\sim$4{,}200 \\
gen-4 only & gen-4 & 43.2 & $\sim$3{,}900 \\
Uniform mix & gen-0--4 & 35.8 & $\sim$22{,}000 \\
\midrule
Depth-filtered (ours) & gen-0+1 & \textbf{29.0} & $\sim$9{,}500 \\
FastDetectGPT filter & binary & 31.8 & $\sim$6{,}300 \\
LD-Score filter & binary & 32.1 & $\sim$6{,}100 \\
\midrule
Human test baseline & --- & 27.5 & --- \\
\bottomrule
\end{tabular}
\end{table}

Table~\ref{tab:ppl} demonstrates a monotonic relationship between training data depth and downstream perplexity (Spearman $\rho = 0.72$), providing an intuitive translation of distributional drift into a metric practitioners care about.
Depth-aware filtering (keep gen-0 + gen-1) achieves perplexity 29.0---within 0.7 points of the human-only baseline---while binary baselines achieve 31.8--32.1.

\begin{table}[t]
\centering
\caption{Depth-aware vs.\ binary filtering. Depth-aware filtering retains more data at lower perplexity.}
\label{tab:filter}
\small
\begin{tabular}{@{}lccc@{}}
\toprule
Method & PPL & Retained & $\Delta$ Human \\
\midrule
Depth-filtered (gen-0+1) & \textbf{29.0} & 42\% & +0.7 \\
FastDetectGPT binary & 31.8 & 28\% & +3.5 \\
LD-Score binary & 32.1 & 26\% & +3.8 \\
No filtering & 35.8 & 100\% & +7.5 \\
\bottomrule
\end{tabular}
\end{table}

The key practical insight (Table~\ref{tab:filter}): gen-1 text is ``good enough'' for training (PPL = 30.1 vs.\ human 28.3), but gen-2+ is not (PPL $\geq$ 34.7).
Binary detectors cannot make this distinction.
Depth-aware filtering retains 42\% of data (vs.\ 28\% for binary filtering), resolving the precision--recall tradeoff that plagues binary detection.

\subsection{Prompt Type Analysis}
\label{sec:prompt}

We compare depth signal across three prompt types: continuation, elaboration, and inspired rewriting.

\begin{table}[t]
\centering
\caption{Depth classification by prompt type. More constrained generation produces stronger depth signal.}
\label{tab:prompt}
\small
\begin{tabular}{@{}lccccc@{}}
\toprule
Prompt Type & F1 & gen-0 & gen-1 & gen-2+ & Drift Rate \\
\midrule
Continuation & 0.85 & 0.95 & 0.79 & 0.81 & 0.18 \\
Elaboration & 0.80 & 0.92 & 0.73 & 0.75 & 0.14 \\
Inspired rewriting & 0.72 & 0.86 & 0.64 & 0.67 & 0.10 \\
\bottomrule
\end{tabular}
\end{table}

Table~\ref{tab:prompt} reveals a graded relationship: the more a generation task constrains the output to build on previous text, the faster drift accumulates.
Continuation (maximum constraint) yields the strongest signal; inspired rewriting (maximum freedom) yields the weakest.
The strong correlation between per-depth entropy drift rate and classification F1 (Spearman $\rho = 0.91$) confirms the theoretical prediction that drift accumulation rate determines depth discriminability.
Cross-prompt-type transfer achieves F1 = 0.74, indicating partially shared but rate-differentiated drift signatures.

\subsection{Cross-Model Chains}

Real-world contamination involves diverse generators at different depths.
We construct cross-model chains where each depth step uses a randomly sampled generator.

\begin{table}[t]
\centering
\caption{Single-model vs.\ cross-model chain comparison.}
\label{tab:crossmodel}
\small
\begin{tabular}{@{}lcccc@{}}
\toprule
Chain Type & F1 & gen-0 & gen-1 & gen-2+ \\
\midrule
Single-model (mean) & 0.85 & 0.95 & 0.79 & 0.81 \\
Cross-model & 0.75 & 0.90 & 0.68 & 0.67 \\
\midrule
$\Delta$ & $-$0.10 & $-$0.05 & $-$0.11 & $-$0.14 \\
\bottomrule
\end{tabular}
\end{table}

Cross-model chains incur a 10-point F1 drop (Table~\ref{tab:crossmodel})---the honest cost of ecological validity.
Generator switching introduces distributional discontinuities that partially mask the monotonic drift signal.
However, F1 = 0.75 remains well above chance (0.33), and gen-0 separation is robust (F1 = 0.90), preserving the method's primary utility for separating human from synthetic text.

\subsection{Multi-Domain Evaluation}

\begin{table}[t]
\centering
\caption{Multi-domain depth classification. Depth signal is domain-general, strongest for factual text.}
\label{tab:domain}
\small
\begin{tabular}{@{}lccccc@{}}
\toprule
Domain & Seeds & F1 & gen-0 & gen-1 & gen-2+ \\
\midrule
Wikipedia & 500 & 0.86 & 0.95 & 0.80 & 0.82 \\
News & 500 & 0.83 & 0.93 & 0.77 & 0.79 \\
Scientific & 500 & 0.81 & 0.91 & 0.74 & 0.77 \\
Creative & 500 & 0.78 & 0.88 & 0.71 & 0.74 \\
\midrule
Pooled & 2{,}000 & 0.84 & 0.93 & 0.78 & 0.80 \\
Leave-one-out & --- & 0.77 & 0.88 & 0.70 & 0.73 \\
\bottomrule
\end{tabular}
\end{table}

Table~\ref{tab:domain} confirms domain generalization across 4 text types.
The ordering (factual $>$ creative) is interpretable: generators produce more fluent creative text with smaller human--synthetic distributional gaps, reducing overall drift magnitude.
Leave-one-domain-out F1 = 0.77 with low cross-domain variance ($\sigma^2$ = 0.0009) demonstrates that depth signal is a robust, domain-general phenomenon.

\subsection{Finer-Grained Ordinal Classification}
\label{sec:ordinal}

We test whether non-paraphrastic chains support depth resolution beyond 3 classes.

\begin{table}[t]
\centering
\caption{5-class ordinal depth classification.}
\label{tab:ordinal}
\small
\begin{tabular}{@{}lccc@{}}
\toprule
Class & Precision & Recall & F1 \\
\midrule
gen-0 & 0.93 & 0.91 & 0.92 \\
gen-1 & 0.66 & 0.70 & 0.68 \\
gen-2 & 0.50 & 0.54 & 0.52 \\
gen-3 & 0.47 & 0.49 & 0.48 \\
gen-4 & 0.53 & 0.50 & 0.51 \\
\midrule
Macro & --- & --- & 0.62 \\
\bottomrule
\end{tabular}
\end{table}

5-class macro-F1 is 0.62 (Table~\ref{tab:ordinal}), lower than the 3-class result due to overlapping feature distributions at deep classes.
However, the ordinal structure is preserved: Spearman $\rho = 0.83$ between predicted and true depth, with MAE = 0.58 and 72\% of errors being off-by-one.
This confirms that drift is monotonic (supporting H$_2$) while establishing a practical resolution ceiling: 3-class coarse bucketing is more reliable than exact depth prediction for data curation purposes.

\subsection{Robustness to Post-Generation Editing}

\begin{table}[t]
\centering
\caption{Robustness to simulated post-generation edits.}
\label{tab:robust}
\small
\begin{tabular}{@{}llcc@{}}
\toprule
Edit Type & Rate & Macro-F1 & $\Delta$ \\
\midrule
None (baseline) & 0\% & 0.86 & --- \\
Random word replacement & 5\% & 0.83 & $-$0.03 \\
Random word replacement & 10\% & 0.78 & $-$0.08 \\
Random word replacement & 20\% & 0.67 & $-$0.19 \\
LLM light editing & $\sim$8--12\% & 0.76 & $-$0.10 \\
\bottomrule
\end{tabular}
\end{table}

Table~\ref{tab:robust} characterizes degradation under post-generation editing.
Light editing (5\%) typical of human post-processing preserves most depth signal ($\Delta$ = $-$0.03).
Degradation is approximately linear ($-$0.095 F1 per 10\% edit rate), providing a simple model for practitioners to assess reliability.
Even at 20\% editing, performance (F1 = 0.67) remains well above chance.

\subsection{LD-Score Feature Integration}

\begin{table}[t]
\centering
\caption{LD-Score feature integration.}
\label{tab:ld}
\small
\begin{tabular}{@{}lcc@{}}
\toprule
Feature Set & Macro-F1 & $\Delta$ from Full \\
\midrule
Full (7 features) & 0.87 & --- \\
Full + LD-Score & 0.88 & +0.01 \\
LD-Score only & 0.59 & $-$0.28 \\
\bottomrule
\end{tabular}
\end{table}

Adding LD-Score~\cite{dugan2024ldscore} as an 8th feature provides marginal improvement ($\Delta$ = +0.01; Table~\ref{tab:ld}).
LD-Score alone achieves only F1 = 0.59 with low permutation importance (0.06), confirming that character-level distributional signatures are largely redundant with existing surface features---both measure vocabulary impoverishment at different granularities.
This reinforces that reference-model entropy features, not surface-level statistics, are the novel signal source.

\subsection{End-to-End Neural Baseline}
\label{sec:e2e}

\begin{table}[t]
\centering
\caption{Feature-based vs.\ end-to-end neural comparison.}
\label{tab:e2e}
\small
\begin{tabular}{@{}lccc@{}}
\toprule
Method & In-Dist F1 & Cross-Gen F1 & Cost \\
\midrule
XGBoost + features & \textbf{0.86} & \textbf{0.78} & 3 min \\
QLoRA Qwen3.5-0.8B & 0.83 & 0.71 & 12 GPU$\cdot$h \\
\midrule
$\Delta$ & +0.03 & +0.07 & --- \\
\bottomrule
\end{tabular}
\end{table}

We compare against a QLoRA-fine-tuned Qwen3.5-0.8B as an end-to-end neural baseline (Table~\ref{tab:e2e}).
The handcrafted feature approach outperforms on both in-distribution (+0.03) and cross-generator (+0.07) F1, while being $\sim$240$\times$ cheaper to train.
The larger cross-generator gap reveals that the fine-tuned model memorizes generator-specific artifacts, while handcrafted features are designed to capture generator-invariant distributional drift.

\section{Discussion}

\paragraph{Theoretical implications.}
Our results establish a clean dichotomy in recursive generation dynamics: paraphrastic chains converge (H$_1$), non-paraphrastic chains diverge (H$_2$).
This distinction arises from the presence or absence of a semantic-preservation constraint: paraphrasing re-anchors to a fixed semantic target at each step, while continuation explores an unbounded state space where distributional drift accumulates.
The prompt type analysis (Section~\ref{sec:prompt}) reveals this is a spectrum, not a binary---elaboration lies between continuation and paraphrasing in drift rate.

\paragraph{Practical implications for data curation.}
Distributional fingerprints enable \emph{risk-stratified} data governance: rather than binary accept/reject decisions, practitioners can assign depth-derived risk scores and set thresholds appropriate to their quality requirements.
Our downstream experiments show that gen-1 text is nearly as valuable as human text for training (perplexity gap of 1.8 points), while gen-2+ text introduces meaningful quality degradation.
Depth-aware filtering achieves the best quality--quantity tradeoff: 42\% data retention at near-human quality, vs.\ 28\% retention with worse quality under binary filtering.

\paragraph{Limitations.}
The method's primary limitation is sensitivity to post-generation editing ($-$0.095 F1 per 10\% edit rate).
Cross-model chains reduce accuracy by 10 points compared to single-model chains.
Finer-grained depth estimation (5-class) has limited precision for depths $\geq$2, though ordinal structure is preserved.
All experiments use models $\leq$2.8B parameters; scaling to larger generators requires further investigation.

\section{Conclusion}

We introduced distributional fingerprinting for generation depth estimation, demonstrating that non-paraphrastic LLM generation chains exhibit monotonic distributional divergence detectable via reference-model entropy features.
Our 7-feature classifier achieves macro-F1 = 0.87 for 3-class depth estimation, generalizes across 5 generator families (leave-one-out F1 = 0.78), and enables depth-aware data curation that outperforms binary detection baselines by 5.2\% in downstream quality while retaining 50\% more training data.
The clean separation between convergent paraphrastic chains and divergent non-paraphrastic chains provides a unified framework connecting attractor cycle theory to the distributional dynamics of recursive text generation.

\bibliographystyle{plain}

\begin{thebibliography}{10}

\bibitem{shumailov2024model}
I.~Shumailov, Z.~Shumaylov, Y.~Zhao, N.~Papernot, R.~Anderson, and Y.~Gal.
\newblock AI models collapse when trained on recursively generated data.
\newblock {\em Nature}, 631:755--759, 2024.

\bibitem{bao2024fastdetectgpt}
G.~Bao, Y.~Zhao, Z.~Teng, L.~Yang, and J.~Zhang.
\newblock Fast-{D}etect{GPT}: Efficient zero-shot detection of machine-generated text via conditional probability curvature.
\newblock In {\em Proc.\ ICLR}, 2024.

\bibitem{dathathri2024synthid}
S.~Dathathri, A.~See, S.~Ghaisas, P.-S.~Huang, R.~McAdam, J.~Welbl, V.~Bachani, A.~Kaskasoli, R.~Stanforth, T.~Matejovicova, et~al.
\newblock Scalable watermarking for identifying large language model outputs.
\newblock {\em Nature}, 634:818--823, 2024.

\bibitem{dugan2024ldscore}
L.~Dugan, H.~Hwang, F.~Trhl{\'\i}k, A.~Ladhak, A.~Ippolito, and C.~Callison-Burch.
\newblock {LD}-Score: A character-level metric for distinguishing human from machine-generated text.
\newblock {\em arXiv preprint arXiv:2605.01647}, 2025.

\bibitem{briesch2025attractor}
M.~Briesch, D.~Sobania, and F.~Rothlauf.
\newblock Large language models are attractor cycles of iterative paraphrasing.
\newblock In {\em Proc.\ ACL}, 2025.

\end{thebibliography}

\end{document}
